\SourceLeaf{196}\Needspace{10\baselineskip}\begin{center}
\textsc{Bemerkung.}
\end{center}

In einem Notizbuch von \textsc{Gauss} steht unter Aufzeichnungen aus den
Jahren 1824 bis 1828 der Vermerk:
\begin{center}»\emph{Journal de l'École Normale. Fouriers
Defin. der Ebene. Lacroix I. p. 503.}«\end{center}

Tatsächlich enthalten die Séances des Écoles Normales T. I (Séance du 25
pluviose an III (1795)), Nouvelle Édition, Paris 1800, S. 28 Äußerungen
\textsc{Fouriers}, in denen er neue Definitionen der Geraden und der Ebene
vorschlägt. \textsc{Fourier} beginnt dabei, ähnlich wie \textsc{Bolyai} und
\textsc{Lobatschefskij}, mit der Kugel. Weitere Veröffentlichungen
\textsc{Fouriers} über die Grundlagen der Geometrie scheinen nicht
zu existieren.

Der Discours préliminaire der Éléments de Géométrie von \textsc{Lacroix}
(Quatrième édition. Paris, An VIII [1804] S. \textsc{xv}.) enthält einen Hinweis auf
das \emph{Journal des Séances de l'École Normale}. Dass dabei
\textsc{Fourier} nicht genannt wird, wohl aber \textsc{Laplace}, erklärt sich
daraus, dass jener seine Theorie in den von diesem geleiteten geometrischen
Übungen der École Normale vorgetragen hatte.

\begin{flushright}
\textsc{Stäckel.}
\end{flushright}

\Needspace{10\baselineskip}
\medskip

\begin{center}
\textsc{Gauss} an \textsc{Schumacher}. Göttingen, 11. Oktober 1827.
\end{center}

.\ .\ .\ .\ .\ .\ .\ .\ .\ Wer ist wohl Herr \textsc{Koch} auf dem Cremon Nr. 83? Er hat
mir seine Parallelentheorie\HistoricalNote{[*)]}{[*)}{Christian Adolf \textsc{Koch}, \emph{Über Parallellinien.} Ein Versuch, dem Urtheil Sachkundiger gewidmet. Hamburg 1827. $8^{0}$. 12 S.]} zugeschickt, an der freilich nichts ist. Aber
es ist etwas Seltenes, dass er den von mir aufgezeigten Mangel sogleich anerkannt hat,
ebenso wie meine Bemerkungen zu seinen späteren, ebenso erfolglosen Versuchen, die er mir schriftlich
zusandte. .\ .\ .\ .\ .

